\chapter{The Nullspace Method}
\label{chap:nullspace}

\vspace*{\fill}

The nullspace method solves the KKT system~\eqref{eq:kkt-system} in two smaller steps. The
search direction is written as $p = Z\eta$, where the columns of $Z$ span the nullspace of
the constraint matrix. This turns the KKT system into a \textbf{reduced system} with the
reduced Hessian, whose dimension equals the number of degrees of freedom of the current face.
The Lagrange multipliers are then recovered from the remaining equations.

\medskip

For the regularized master problem both steps can be carried out explicitly. The reordering
theorem of the previous chapter gives a nullspace basis that consists of the vectors
$e_j - e_{\rep(j)}$ and of unit vectors for the free surpluses. This basis is
\textbf{never stored}: it is determined by the five sets alone.

\medskip

The central result of this chapter is a \textbf{decoupling} of the reduced system. The
surplus part of the direction is given by an explicit formula, and only an
$|\Nx|\times|\Nx|$ system has to be solved. Its matrix is the \textbf{effective reduced
Hessian} $\Heff = (\Aeff Z_x)^\top(\Aeff Z_x)$, the reduced Hessian of the equality rows
together with the inequality rows whose surplus is active. Finally, the multipliers turn out
to be the reduced costs of the columns and the duals of the inequality rows, available at
almost no cost.

\vspace*{\fill}

\newpage


\section{The Nullspace Basis}
\label{sec:basis}

By \cref{thm:rank}, the nullspace of $D_\Wset$ has dimension $|\Nx|+|\Ns|$. We label its basis
vectors by the indices of these two sets.

\begin{theorem}[Nullspace basis] \label{thm:basis}
Under the assumptions of \cref{thm:reordering}, the matrix
\begin{equation} \label{eq:Z-reordered}
Z = Q\tilde Z,\qquad
\tilde Z = \begin{bmatrix}
-C_{\Nx} & 0\\
0 & 0\\
I_{|\Nx|} & 0\\
0 & 0\\
0 & I_{|\Ns|}
\end{bmatrix}
\begin{matrix} \leftarrow\Sset\\ \leftarrow\Wx\\ \leftarrow\Nx\\ \leftarrow\Ws\\ \leftarrow\Ns\end{matrix}
\end{equation}
is a basis of $\Null(D_\Wset)$. In the original order of the variables it is block-diagonal,
\begin{equation} \label{eq:Z-blocks}
Z = \begin{bmatrix} Z_x & 0\\ 0 & Z_s\end{bmatrix},
\qquad Z_x\in\R^{n\times|\Nx|},\quad Z_s\in\R^{m_I\times|\Ns|},
\end{equation}
where $Z_x$ is a nullspace basis for the column weights and $Z_s$ is the selection matrix of
the free surpluses.
\end{theorem}

\begin{proof}
Write $\tilde Z = Q^\top Z$ in the five row blocks of the extended ordering,
\[
\tilde Z = \bigl[\,Z_\Sset;\; Z_{\Wx};\; Z_{\Nx};\; Z_{\Ws};\; Z_{\Ns}\,\bigr].
\]
By~\eqref{eq:reordered},
$D_\Wset Z = (D_\Wset Q)\tilde Z = 0$ is equivalent to
\[
Z_\Sset + C_{\Wx}Z_{\Wx} + C_{\Nx}Z_{\Nx} = 0,\qquad Z_{\Wx} = 0,\qquad Z_{\Ws} = 0 .
\]
The blocks $Z_{\Nx}$ and $Z_{\Ns}$ are free, so the nullspace has the stated dimension, and the
choice $Z_{\Nx} = I$, $Z_{\Ns} = I$ (with the corresponding zero columns) gives a basis with
$Z_\Sset = -C_{\Nx}$ in the first group of columns and $Z_\Sset = 0$ in the second. The column
weights and the surpluses do not mix, which gives~\eqref{eq:Z-blocks}. 
\end{proof}

\begin{cor}[Explicit nullspace vectors] \label{cor:z-explicit}
The columns of $Z$ are
\begin{equation} \label{eq:z-explicit}
z_j = e_j - e_{\rep(j)}\quad (j\in\Nx),\qquad\qquad z_i = e_{n+i}\quad (i\in\Ns),
\end{equation}
where $e_\ell$ denotes the $\ell$-th unit vector in $\R^{n+m_I}$.
\end{cor}

\begin{proof}
Column $j\in\Nx$ of $\tilde Z$ has a $1$ in the position of $j$ and the entries of
$-C_{\Nx}e_j = -e_{\psi(j)}$ in the representative block. By the column generation property,
this is a single $-1$ in the row of block $\psi(j)$, i.e.\ at the position of the
representative $S_{\psi(j)} = \rep(j)$. Column $i\in\Ns$ has a single $1$ at the position of
surplus $i$. 
\end{proof}

\begin{figure}[ht]
\centering
\begin{tikzpicture}[x=0.75cm, y=1cm, font=\small]
    % one block with 5 columns: rep = 2, N = {1,4}, W = {3,5}
    \foreach \j/\h/\lab/\col in {1/0.25/j_1/cblue, 2/0.5/{\rep}/cgreen, 3/0/w/cgrey, 4/0.25/j_2/cblue, 5/0/w/cgrey} {
        \fill[\col!45] (\j-0.35,0) rectangle (\j+0.35,{\h*3});
        \draw[\col!80!black] (\j-0.35,0) rectangle (\j+0.35,{\h*3});
    }
    \node[below] at (1,0) {$x_{j_1}$};
    \node[below] at (2,0) {$x_{\rep}$};
    \node[below] at (3,0) {$x_{w}$};
    \node[below] at (4,0) {$x_{j_2}$};
    \node[below] at (5,0) {$x_{w'}$};
    \draw (0.4,0) -- (5.6,0);
    % arrows for z_{j1}: +1 at j1, -1 at rep
    \draw[-{Stealth[length=2.5mm]}, thick, cred] (1,0.85) -- (1,1.45) node[above] {$+\eta_{j_1}$};
    \draw[-{Stealth[length=2.5mm]}, thick, cred] (2,1.6) -- (2,1.0) node[pos=0, above] {$-\eta_{j_1}$};
    \node[align=left, right] at (6.2,1.1) {$x + \eta_{j_1} z_{j_1}$:\\ weight moves from the\\ representative to $j_1$;\\ the block sum stays $1$,\\ and $x_w = x_{w'} = 0$ stay.};
    % surplus
    \begin{scope}[xshift=11.5cm]
        \fill[corange!45] (-0.35,0) rectangle (0.35,0.9);
        \draw[corange!80!black] (-0.35,0) rectangle (0.35,0.9);
        \node[below] at (0,0) {$s_i$};
        \draw (-0.6,0) -- (0.6,0);
        \draw[-{Stealth[length=2.5mm]}, thick, cred] (0,1.0) -- (0,1.6) node[above] {$+\eta_i$};
        \node[align=left, right] at (0.8,0.8) {$z_i = e_{n+i}$:\\ a free surplus\\ moves on its own};
    \end{scope}
\end{tikzpicture}
\caption{Interpretation of the nullspace vectors~\eqref{eq:z-explicit}. A vector
$z_j = e_j - e_{\rep(j)}$ transfers weight between a nullspace column and the
representative of its block, which keeps the convexity constraint satisfied and the
working-set columns at zero. A free surplus belongs to no block and moves independently.}
\label{fig:nullspace-vectors}
\end{figure}

\Cref{cor:z-explicit} shows that $Z$ is completely determined by the extended ordering: by
the sets $\Nx$, $\Ns$ and the representative mapping. It therefore never has to be formed.
The same holds for the product of the linking matrix with the nullspace vectors,
\begin{equation} \label{eq:AZ-columns}
\Abar z_j = \begin{bmatrix} a_j^E - a_{\rep(j)}^E\\ a_j^I - a_{\rep(j)}^I\end{bmatrix}
= a_j - a_{\rep(j)} =: u_j\quad(j\in\Nx),\qquad
\Abar z_i = -e_i^{(I)}\quad(i\in\Ns),
\end{equation}
where $e_i^{(I)}\in\R^m$ is the unit vector of inequality row $i$.


\section{The Reduced KKT System}
\label{sec:reduced}

\subsection{The reduced system}

Every direction that satisfies the constraints of~\eqref{eq:eqp} has the form
$p = Z\eta$ with $\eta = [\eta_x;\eta_s]\in\R^{|\Nx|+|\Ns|}$. Substituting this into the first
block row of~\eqref{eq:kkt-system} and multiplying by $Z^\top$, which annihilates
$D_\Wset^\top$, gives the \textbf{reduced KKT system}
\begin{equation} \label{eq:reduced-system}
\frac{1}{\delta}\,Z^\top\Abar^\top\Abar Z\,\eta = -Z^\top g .
\end{equation}
We call $\Hbar = Z^\top\Abar^\top\Abar Z$ the (unscaled) \textbf{reduced Hessian}. The vector
$\eta$ minimizes $f(y_k + Z\eta) = f(y_k) + (Z^\top g)^\top\eta + \frac{1}{2\delta}\eta^\top\Hbar\eta$,
so~\eqref{eq:reduced-system} is the subproblem~\eqref{eq:eqp} expressed in the coordinates
of the nullspace. The parameter $\delta$ only appears on the right-hand side if we write the
system as $\Hbar\eta = -\delta Z^\top g$.

\subsection{Block structure}

To exhibit the structure of $\Hbar$, we introduce the products of the linking matrices with
the column part of the basis,
\[
B_E := A_E Z_x\in\R^{m_E\times|\Nx|},\qquad B_I := A_I Z_x\in\R^{m_I\times|\Nx|},
\]
and split the rows of $B_I$ and of $\pi_I$ according to the status of the surplus of each row:
\[
B_\Wset := W_s B_I,\quad \pi_\Wset := W_s\pi_I \qquad\text{(rows with an active surplus)},
\]
\[
B_\mathcal{N} := Z_s^\top B_I,\quad \pi_\mathcal{N} := Z_s^\top\pi_I \qquad\text{(rows with a free surplus)}.
\]
Since every inequality row belongs to exactly one of the two groups, $W_s^\top W_s + Z_sZ_s^\top = I$
and therefore
\begin{equation} \label{eq:BI-split}
B_I^\top B_I = B_\Wset^\top B_\Wset + B_\mathcal{N}^\top B_\mathcal{N} .
\end{equation}

\begin{lemma}[Blocks of the reduced system] \label{lem:reduced-blocks}
With~\eqref{eq:Z-blocks},
\[
\Abar Z = \begin{bmatrix} B_E & 0\\ B_I & -Z_s\end{bmatrix},\qquad
\Hbar = \begin{bmatrix} H & -B_\mathcal{N}^\top\\ -B_\mathcal{N} & I\end{bmatrix},\qquad
Z^\top g = \begin{bmatrix} Z_x^\top g_x\\ \pi_\mathcal{N}\end{bmatrix},
\]
where $H = B_E^\top B_E + B_I^\top B_I = (AZ_x)^\top(AZ_x)$ is the reduced Hessian of the column
weights with all linking rows.
\end{lemma}

\begin{proof}
The first identity follows from the block forms of $\Abar$ and $Z$. Then
$\Hbar = (\Abar Z)^\top(\Abar Z)$ has the blocks $B_E^\top B_E + B_I^\top B_I$,
$-B_I^\top Z_s = -B_\mathcal{N}^\top$ and $Z_s^\top Z_s = I$. Finally
$Z^\top g = [Z_x^\top g_x;\, Z_s^\top g_s]$ and $g_s = \pi_I$ by~\eqref{eq:gradient-split}.

\end{proof}

\subsection{Decoupling and the effective reduced Hessian}

The identity block in the lower right corner of $\Hbar$ allows us to eliminate the surplus
part of the direction.

\begin{theorem}[Decoupling of the reduced system] \label{thm:decoupling}
The reduced system~\eqref{eq:reduced-system} is equivalent to
\begin{align}
\Heff\,\eta_x &= -\delta\,Z_x^\top\bigl(c - A_E^\top\pi_E - A_I^\top W_s^\top\pi_\Wset\bigr),
\label{eq:eff-system}\\[3pt]
\eta_s &= B_\mathcal{N}\,\eta_x - \delta\,\pi_\mathcal{N},
\label{eq:surplus-direction}
\end{align}
with the \textbf{effective reduced Hessian}
\begin{equation} \label{eq:Heff}
\Heff := H - B_\mathcal{N}^\top B_\mathcal{N} = B_E^\top B_E + B_\Wset^\top B_\Wset .
\end{equation}
\end{theorem}

\begin{proof}
Multiply~\eqref{eq:reduced-system} by $\delta$ and write it in the blocks of
\cref{lem:reduced-blocks}:
\[
H\eta_x - B_\mathcal{N}^\top\eta_s = -\delta Z_x^\top g_x,\qquad
-B_\mathcal{N}\eta_x + \eta_s = -\delta\pi_\mathcal{N}.
\]
The second equation is~\eqref{eq:surplus-direction}. Substituting it into the first gives
$(H - B_\mathcal{N}^\top B_\mathcal{N})\eta_x = -\delta\bigl(Z_x^\top g_x + B_\mathcal{N}^\top\pi_\mathcal{N}\bigr)$.
The second expression for $\Heff$ follows from~\eqref{eq:BI-split}. For the right-hand side,
insert $g_x = c - A_E^\top\pi_E - A_I^\top\pi_I$ and split
$Z_x^\top A_I^\top\pi_I = B_I^\top\pi_I = B_\Wset^\top\pi_\Wset + B_\mathcal{N}^\top\pi_\mathcal{N}$.
The terms with $\pi_\mathcal{N}$ cancel, and $B_\Wset^\top\pi_\Wset = Z_x^\top A_I^\top W_s^\top\pi_\Wset$.

\end{proof}

\begin{figure}[ht]
\centering
\begin{tikzpicture}[x=0.5cm, y=0.5cm, font=\small]
    % H bar
    \fill[cblue!25] (0,3) rectangle (6,9);
    \fill[corange!30] (6,3) rectangle (9,9);
    \fill[corange!30] (0,0) rectangle (6,3);
    \fill[cgreen!25] (6,0) rectangle (9,3);
    \draw[thick] (0,0) rectangle (9,9);
    \draw (6,0) -- (6,9); \draw (0,3) -- (9,3);
    \draw[cgreen!60!black, very thick] (6,3) -- (9,0);
    \node at (3,6) {$H$};
    \node at (7.5,6) {$-B_\mathcal{N}^\top$};
    \node at (3,1.5) {$-B_\mathcal{N}$};
    \node[fill=cgreen!25, inner sep=1pt] at (7.9,2.1) {$I$};
    \node at (4.5,10) {$\Hbar$};
    \node[left] at (0,6) {$\Nx$}; \node[left] at (0,1.5) {$\Ns$};
    % arrow
    \draw[-{Stealth[length=3mm]}, thick] (10,4.5) -- (14,4.5)
        node[midway, above] {eliminate $\eta_s$};
    % H eff
    \begin{scope}[xshift=7.5cm]
    \fill[cblue!25] (0,1.5) rectangle (6,7.5);
    \draw[thick] (0,1.5) rectangle (6,7.5);
    \node[align=center] at (3,4.5) {$\Heff =$\\ $H - B_\mathcal{N}^\top B_\mathcal{N}$};
    \node at (3,8.5) {$|\Nx|\times|\Nx|$};
    \end{scope}
\end{tikzpicture}
\caption{The reduced Hessian $\Hbar$ of \cref{lem:reduced-blocks} and its Schur complement.
Eliminating the identity block of the free surpluses leaves the effective reduced Hessian
$\Heff$ on the nullspace columns only.}
\label{fig:schur}
\end{figure}

$\Heff$ is the Schur complement of the identity block in $\Hbar$
(\cref{fig:schur}). Only an $|\Nx|\times|\Nx|$ system has to be solved; the surplus
part~\eqref{eq:surplus-direction} costs one sparse matrix--vector product. The following
corollary gives the interpretation that the rest of this work builds on.

\begin{cor}[Effective linking matrix] \label{cor:Aeff}
Define the \textbf{effective linking matrix} and the corresponding duals
\[
\Aeff = \begin{bmatrix} A_E\\ W_sA_I\end{bmatrix},\qquad
\pieff = \begin{bmatrix}\pi_E\\ \pi_\Wset\end{bmatrix},
\]
i.e.\ the equality rows together with the inequality rows whose surplus is active. Then
\begin{equation} \label{eq:Heff-Aeff}
\Heff = (\Aeff Z_x)^\top(\Aeff Z_x),\qquad
\Heff\,\eta_x = -\delta\,Z_x^\top\bigl(c - \Aeff^\top\pieff\bigr).
\end{equation}
This is the reduced system of the master problem with equality linking rows only, whose
linking rows are the rows of $\Aeff$.
\end{cor}

\begin{proof}
$\Aeff Z_x = [B_E;\, B_\Wset]$, which gives the first identity by~\eqref{eq:Heff}. The second is
the right-hand side of~\eqref{eq:eff-system}, since
$\Aeff^\top\pieff = A_E^\top\pi_E + A_I^\top W_s^\top\pi_\Wset$. 
\end{proof}

\begin{remark}[Interpretation] \label{rem:Aeff-interpretation}
An inequality row with a \emph{free} surplus does not restrict the column weights: any
change of its residual is absorbed by the surplus, so the row contributes neither curvature
nor dual pressure to the column part of the problem. An inequality row with an \emph{active}
surplus behaves like an equality row. The column part of the method is therefore the
method for equality rows, applied to a linking matrix $\Aeff$ whose rows change with the
surplus working set. Without inequality rows, or with all surpluses active, $\Aeff = A$ and
$\Heff = H$.
\end{remark}

\begin{cor}[Positive definiteness and dimension] \label{cor:dimension}
$\Hbar$ is positive definite if and only if $\Heff$ is, i.e.\ if and only if $\Aeff Z_x$ has
full column rank. In this case
\[
|\Nx|\;\le\;\rank(\Aeff)\;\le\;m_E + |\Ws|,\qquad\text{and hence}\qquad |\Nx| + |\Ns|\le m .
\]
\end{cor}

\begin{proof}
A symmetric block matrix with a positive definite diagonal block is positive definite if
and only if the Schur complement of that block is (Haynsworth inertia additivity). The
rank bound holds because $\Heff = (\Aeff Z_x)^\top(\Aeff Z_x)$ is nonsingular only if
$|\Nx|\le\rank(\Aeff Z_x)\le\rank(\Aeff)$, and $\Aeff$ has $m_E+|\Ws|$ rows. Adding
$|\Ns| = m_I - |\Ws|$ gives the last bound. 
\end{proof}

\Cref{cor:dimension} is the reason the method scales: the system that has to be factorized
never exceeds the number of linking rows, \emph{however many columns are generated}.

\subsection{The search direction}

With $\eta$ from \cref{thm:decoupling}, the direction $p = Z\eta$ has the components
\begin{equation} \label{eq:direction-components}
p_j = \eta_j\;(j\in\Nx),\qquad
p_{S_k} = -\!\!\sum_{j\in\Nx,\,\psi(j) = k}\!\!\eta_j\;(k = 1,\dots,K),\qquad
(p_s)_i = (A_Ip_x)_i - \delta\,\pi_i\;(i\in\Ns),
\end{equation}
and zero on $\Wx$ and $\Ws$. The surplus component uses $B_\mathcal{N}\eta_x = Z_s^\top A_I Z_x\eta_x = Z_s^\top A_I p_x$.
The direction is supported on $\Sset\cup\Nx\cup\Ns$, and its column part and the vector
$A_Ip_x$ only involve the columns in $\Sset\cup\Nx$.

\begin{prop}[Full step] \label{prop:full-step}
If the ratio test returns $\alpha = 1$, the new iterate $y_k+p$ is a subspace minimizer, i.e.\
the next reduced gradient $Z^\top\nabla f(y_k+p)$ vanishes.
\end{prop}

\begin{proof}
On the affine set $\{y_k + Z\eta\}$, $f$ is the convex quadratic
$\phi(\eta) = f(y_k) + (Z^\top g)^\top\eta + \frac{1}{2\delta}\eta^\top\Hbar\eta$, whose
minimizer is the solution of~\eqref{eq:reduced-system}. 
\end{proof}

After a full step the method can therefore skip the next direction computation and test the
multipliers directly. The next proposition shows how the free surpluses act on the duals.

\begin{prop}[Duals of free-surplus rows] \label{prop:free-duals}
A step $y\gets y + \alpha p$ changes the dual of every row with a free surplus as
\[
\pi_i\;\gets\;(1-\alpha)\,\pi_i\qquad(i\in\Ns).
\]
In particular $\pi_i = 0$ on these rows at every subspace minimizer.
\end{prop}

\begin{proof}
By~\eqref{eq:pi-of-y}, $\pi$ changes by $-\frac{\alpha}{\delta}\Abar p$. For row $i\in\Ns$,
$(\Abar p)_i = (A_Ip_x)_i - (p_s)_i = \delta\pi_i$ by~\eqref{eq:direction-components}. 
\end{proof}

\begin{remark}[Independence of $\delta$ and $\pihat$] \label{rem:factor-independent}
$\Heff$ depends only on the columns and the working set, not on $\delta$, $\pihat$ or the
iterate. A change of the stability center or of $\delta$ only changes the right-hand side
of~\eqref{eq:eff-system}. This is what makes the hot start of \cref{sec:hot-start} possible.
\end{remark}

\subsection{Implicit evaluation}

Neither $Z$, nor $\Aeff Z_x$, nor the gradient has to be stored. With $u_j = a_j - a_{\rep(j)}$
from~\eqref{eq:AZ-columns} and the diagonal \textbf{row mask} $D_\rho\in\R^{m\times m}$ with
$\rho_\ell = 0$ for the rows of free surpluses and $\rho_\ell = 1$ otherwise:
\begin{equation} \label{eq:implicit}
(\Heff)_{kj} = u_k^\top D_\rho\,u_j,\qquad
\bigl(Z_x^\top(c - \Aeff^\top\pieff)\bigr)_j = \bigl(c_j - a_j^\top D_\rho\pi\bigr)
- \bigl(c_{\rep(j)} - a_{\rep(j)}^\top D_\rho\pi\bigr).
\end{equation}
Both are computed from the sparse columns $a_j$ and the $m$-vector $\pi$. The cost of one entry
of the right-hand side is proportional to the number of nonzeros of two columns, independent
of the total number of columns $n$.


\section{The Lagrange Multipliers}
\label{sec:multipliers}

The multipliers follow from the first block row of the KKT system~\eqref{eq:kkt-system}:
\begin{equation} \label{eq:multiplier-system}
D_\Wset^\top\begin{bmatrix}\mu\\ \nu_x\\ \nu_s\end{bmatrix} = h,\qquad
h := \frac{1}{\delta}\Abar^\top\Abar\,p + g .
\end{equation}

\begin{theorem}[Lagrange multipliers] \label{thm:multipliers}
With the extended ordering of \cref{thm:reordering} and $h = [h_\Sset; h_{\Wx}; h_{\Nx}; h_{\Ws}; h_{\Ns}]$,
the solution of~\eqref{eq:multiplier-system} is
\begin{equation} \label{eq:multipliers}
\mu = h_\Sset,\qquad \nu_x = h_{\Wx} - C_{\Wx}^\top\mu,\qquad \nu_s = h_{\Ws},
\end{equation}
and moreover $h_{\Nx} = C_{\Nx}^\top\mu$ and $h_{\Ns} = 0$.
\end{theorem}

\begin{proof}
Multiply~\eqref{eq:multiplier-system} by $Q^\top$ and use the transpose of~\eqref{eq:reordered}:
\[
\begin{bmatrix}
I_K & 0 & 0\\
C_{\Wx}^\top & I & 0\\
C_{\Nx}^\top & 0 & 0\\
0 & 0 & I\\
0 & 0 & 0
\end{bmatrix}
\begin{bmatrix}\mu\\ \nu_x\\ \nu_s\end{bmatrix}
=
\begin{bmatrix} h_\Sset\\ h_{\Wx}\\ h_{\Nx}\\ h_{\Ws}\\ h_{\Ns}\end{bmatrix}.
\]
The first, second and fourth block rows give~\eqref{eq:multipliers}; the third and fifth
are the remaining identities. 
\end{proof}

The multipliers are only needed at a subspace minimizer, where $p = 0$ and therefore $h = g$.
In this case \cref{thm:multipliers} takes a particularly simple form.

\begin{cor}[Multipliers at a subspace minimizer] \label{cor:multipliers-minimizer}
At a subspace minimizer, with $\pi = \pi(y)$,
\begin{align*}
\mu_k &= c_{S_k} - a_{S_k}^\top\pi && (k = 1,\dots,K),\\
(\nu_x)_j &= c_j - a_j^\top\pi - \mu_{\psi(j)} = \bar c_j(\pi,\mu) && (j\in\Wx),\\
(\nu_s)_i &= \pi_i && (i\in\Ws),
\end{align*}
and $\bar c_j(\pi,\mu) = 0$ for $j\in\Sset\cup\Nx$, $\pi_i = 0$ for $i\in\Ns$.
\end{cor}

\begin{proof}
Insert $h = g$ and~\eqref{eq:gradient-split} into \cref{thm:multipliers}. For $j\in\Nx$,
$h_{\Nx} = C_{\Nx}^\top\mu$ reads $g_j = \mu_{\psi(j)}$, i.e.\ $\bar c_j = 0$, and $h_{\Ns} = 0$
reads $\pi_i = 0$. 
\end{proof}

\Cref{cor:multipliers-minimizer} gives a concrete form of \cref{rem:same-test}: the optimality
test of the active-set method checks the reduced costs of the working-set columns and the
signs of the duals of the inequality rows with active surpluses. All of them are computed
from the single vector $\pi$. The test touches every column of $\Wx$ once and costs
$O(\nnz(A))$; it is the only operation of the method whose cost grows with the number of
generated columns.

\begin{remark}[``$\le$'' rows] \label{rem:kappa-summary}
For a ``$\le$'' row with slack coefficient $\kappa_i = +1$ (\cref{rem:kappa}), the formulas
above change as follows; everything else, in particular $\Heff$ and $\Aeff$, is unchanged
because $\kappa_i^2 = 1$:
\[
\pi_I = \pihat_I + \tfrac{1}{\delta}\bigl(b_I - A_Ix - K_s s\bigr),\qquad
g_{n+i} = -\kappa_i\pi_i,\qquad
(\nu_s)_i = -\kappa_i\pi_i,\qquad
(p_s)_i = \kappa_i\bigl(\delta\pi_i - (A_Ip_x)_i\bigr).
\]
For $\kappa_i = -1$ these are the formulas of the text.
\end{remark}


\section{Summary of One Iteration}
\label{sec:iteration-summary}

\Cref{alg:nullspace-iteration} collects the computations of one iteration of the active-set
method in nullspace form, and \cref{tab:costs} lists their costs. Here $\bar\nu$ denotes the
average number of nonzeros of a column $a_j$.

\begin{algobox}{One Iteration of the Active-Set Method in Nullspace Form}{alg:nullspace-iteration}
\State \textbf{Given:} iterate $y = [x;s]$, the extended ordering, $\pi = \pi(y)$, the factor $\Heff = LDL^\top$
\State Compute $\mu^{\rho}_k = c_{S_k} - a_{S_k}^\top D_\rho\pi$ and the right-hand side
       $r_j = -\delta\bigl(c_j - a_j^\top D_\rho\pi - \mu^{\rho}_{\psi(j)}\bigr)$, $j\in\Nx$
\State Solve $LDL^\top\eta_x = r$
\State Expand $p_x = Z_x\eta_x$ and compute $(p_s)_i = (A_Ip_x)_i - \delta\pi_i$, $i\in\Ns$
\If{$p\approx0$}
    \State $\mu_k = c_{S_k} - a_{S_k}^\top\pi$; \;$\nu_x = $ reduced costs on $\Wx$; \;$\nu_s = \pi$ on $\Ws$
    \State Stop if all multipliers are $\ge-\varepsilon$; otherwise reduce the working set
           (\cref{sec:reducing})
\Else
    \State Ratio test over $\Sset\cup\Nx\cup\Ns$; \;$y\gets y+\alpha p$; \;$\pi\gets\pi - \frac{\alpha}{\delta}\Abar p$
    \State If blocked, expand the working set (\cref{sec:expanding})
\EndIf
\end{algobox}

\begin{table}[ht]
\centering
\renewcommand{\arraystretch}{1.25}
\begin{tabular}{lll}
\toprule
operation & cost & grows with $n$? \\
\midrule
right-hand side of~\eqref{eq:eff-system} & $O\bigl((|\Nx|+K)\,\bar\nu\bigr)$ & no \\
solve with the factor of $\Heff$ & $O(\nnz(L))$, at most $O(|\Nx|^2)$ & no \\
surplus direction~\eqref{eq:surplus-direction} & $O\bigl((|\Nx|+K)\,\bar\nu + m_I\bigr)$ & no \\
ratio test, step and update of $\pi$ & $O\bigl((|\Nx|+K)\,\bar\nu + m\bigr)$ & no \\
multipliers (\cref{cor:multipliers-minimizer}) & $O(\nnz(A) + m_I)$ & yes \\
\bottomrule
\end{tabular}
\caption{Cost of the operations of one iteration. Only the multiplier test touches all
columns.}
\label{tab:costs}
\end{table}
